\chapter{Theoretical Background}

This chapter presents the foundational knowledge required to fully understand the
methods described in subsequent chapters. It covers the core concepts of
reinforcement learning, the principal algorithm families employed in this thesis,
the multi-agent extension of the RL framework, and the simulation tools and dataset
used to construct the experimental environment.

% ================================================================
\section{Reinforcement Learning}
% ================================================================

\subsection{Core Concepts and the Markov Decision Process}

Reinforcement Learning is a machine learning paradigm in which an agent learns
to make sequential decisions by interacting directly with an
environment~\cite{sutton2018reinforcement}. At each discrete time step $t$, the
agent observes the current state $s_t \in \mathcal{S}$, selects an action
$a_t \in \mathcal{A}$, receives a scalar reward $r_t$, and transitions to a new
state $s_{t+1}$. The agent's goal is to find a policy $\pi$ that maximises the
expected discounted cumulative return:
\[
  G_t = \sum_{k=0}^{\infty} \gamma^k\, r_{t+k+1},
\]
where $\gamma \in [0,1)$ is a discount factor that controls the relative importance
of future rewards.

This sequential decision-making problem is formally modelled as a
\textit{Markov Decision Process}, defined by the tuple
$(\mathcal{S}, \mathcal{A}, P, R, \gamma)$, where $P$ is the state-transition
function and $R$ is the reward function~\cite{sutton2018reinforcement}. A central
quantity in RL is the action-value function:
\[
  Q^\pi(s,a) = \mathbb{E}_\pi\!\left[G_t \mid s_t = s,\; a_t = a\right],
\]
which represents the expected return when action $a$ is taken in state $s$ and the
agent subsequently follows policy $\pi$. The optimal policy $\pi^*$ is the one that
maximises $Q^\pi(s,a)$ for all state-action pairs.

In the traffic signal control setting, the state describes the traffic condition at
an intersection (e.g.\ queue lengths and vehicle counts per lane), the action
corresponds to the selection of the next signal phase, and the reward reflects the
degree of congestion reduction achieved after each decision.

% ----------------------------------------------------------------
\subsection{Algorithm Taxonomy}
% ----------------------------------------------------------------

RL algorithms are commonly grouped into three families based on what they
learn~\cite{sutton2018reinforcement}.

\textbf{Value-based methods} estimate the action-value function $Q(s,a)$ and derive
a policy by selecting the action with the highest estimated value. The canonical
update rule is:
\[
  Q(s_t,a_t) \leftarrow Q(s_t,a_t)
  + \alpha\!\left[r_t + \gamma\max_{a'} Q(s_{t+1},a') - Q(s_t,a_t)\right],
\]
where $\alpha$ is the learning rate. These methods are well suited to discrete
action spaces and form the basis of DQN and DDQN.

\textbf{Policy-based methods} directly optimise a parameterised stochastic policy
$\pi_\theta(a \mid s)$ by computing the gradient of the expected return with
respect to $\theta$. They are not restricted to discrete action spaces and often
exhibit greater stability in high-dimensional state spaces. PPO belongs to this
family.

\textbf{Actor-critic methods} combine both approaches: an \textit{actor} maintains
the policy and selects actions, while a \textit{critic} estimates a value function
to evaluate action quality and guide policy updates. PPO is also an actor-critic
algorithm, as it employs a separate critic network alongside the policy network.

In traffic signal control, the state space grows rapidly with the number of lanes
and intersections, rendering tabular Q-learning intractable. \textit{Deep
Reinforcement Learning} addresses this by using deep neural networks to
approximate the value function or the policy, enabling generalisation over large
and continuous state spaces~\cite{mnih2015humanlevel}. All three algorithms used in
this thesis---DQN, DDQN, and PPO---are DRL methods.

% ================================================================
\section{Deep Q-Network and Double DQN}
% ================================================================

\subsection{Deep Q-Network}

Deep Q-Network, proposed by Mnih et al.\ in 2015, replaces the tabular
Q-function with a deep neural network parameterised by $\theta$, so that
$Q(s,a;\theta)$ can be evaluated for large state spaces~\cite{mnih2015humanlevel}.
The network is trained by minimising the mean-squared Bellman error against a
target:
\[
  y_t = r_t + \gamma\max_{a'} Q(s_{t+1}, a';\, \theta^{-}),
\]
where $\theta^{-}$ denotes the parameters of a separate \textit{target network}
that is periodically copied from the main network. Freezing the target for several
updates stabilises training by preventing the optimisation target from shifting at
every gradient step.

A second key mechanism is \textit{experience replay}: transitions
$(s_t, a_t, r_t, s_{t+1})$ are stored in a replay buffer and sampled uniformly at
random during training. This breaks the temporal correlation between consecutive
samples and improves data efficiency~\cite{mnih2015humanlevel}.

\subsection{Double DQN}

A known pathology of DQN is systematic \textit{overestimation} of Q-values, which
arises because the same network is used both to select the greedy action and to
evaluate its value. Double DQN, introduced by van Hasselt et
al.~\cite{hasselt2016ddqn}, decouples these two roles: the main network selects the
action, and the target network evaluates it:
\[
  y_t = r_t + \gamma\, Q\!\left(s_{t+1},\;
        \underset{a'}{\arg\max}\; Q(s_{t+1}, a';\,\theta);\;
        \theta^{-}\right).
\]
This simple modification reduces estimation bias and improves training stability
without adding meaningful computational overhead. In this project, both DQN and
DDQN are implemented via Stable-Baselines3~\cite{raffin2021sb3}, with DDQN
obtained by modifying the target computation of the DQN implementation.

% ================================================================
\section{Proximal Policy Optimization}
% ================================================================

Proximal Policy Optimization, introduced by Schulman et
al.~\cite{schulman2017ppo}, is an actor-critic, policy-gradient algorithm designed
to achieve stable and sample-efficient updates. Rather than learning a value
function and deriving a policy from it, PPO directly optimises the parameterised
stochastic policy $\pi_\theta(a \mid s)$.

The core challenge of policy-gradient methods is that large parameter updates can
cause the policy to change drastically, destabilising learning. PPO addresses this
with a \textit{clipped surrogate objective} that constrains how much the policy may
change in a single update:
\[
  L^{\text{CLIP}}(\theta) = \mathbb{E}_t\!\left[
    \min\!\left(
      r_t(\theta)\,\hat{A}_t,\;
      \operatorname{clip}\!\left(r_t(\theta), 1-\varepsilon, 1+\varepsilon\right)
      \hat{A}_t
    \right)
  \right],
\]
where $r_t(\theta) = \pi_\theta(a_t \mid s_t)/\pi_{\theta_{\text{old}}}(a_t \mid
s_t)$ is the probability ratio between the new and old policies, $\varepsilon$ is
the clip threshold, and $\hat{A}_t$ is an estimate of the advantage function that
measures how much better action $a_t$ is compared to the average action at state
$s_t$.

The advantage is estimated using \textit{Generalised Advantage Estimation}~\cite{schulman2017ppo}, which interpolates between high-variance Monte Carlo
returns and low-variance (but biased) one-step TD estimates, offering a practical
bias-variance trade-off. PPO collects experience in fixed-length rollouts and
performs multiple gradient epochs over each rollout before discarding it---a
fundamentally different data-flow from DQN's replay buffer. PPO is implemented in
this project via Stable-Baselines3~\cite{raffin2021sb3}.

% ================================================================
\section{Multi-Agent Reinforcement Learning}
% ================================================================

\subsection{Problem Formulation}

When the problem is extended from a single intersection to a road network,
each intersection becomes an independent decision-making agent. This gives rise to
\textit{Multi-Agent Reinforcement Learning}, in which multiple agents
interact with a shared environment and must collectively optimise network-wide
traffic performance~\cite{zhang2019integrating}.

The problem is naturally formulated as a \textit{Decentralised Partially Observable
MDP}~\cite{oliehoek2016concise}. Unlike a standard MDP, each agent $i$
receives only a local observation $o_i$ of the global state, selects a local action
$a_i$, and receives a (possibly shared) reward $r_i$. The joint action of all agents
drives the environment's state transition.

\subsection{Centralised vs.\ Decentralised Approaches}

Two broad paradigms exist for solving MARL problems~\cite{terry2020revisiting}.
\textit{Centralised} methods treat the entire network as a single large agent that
observes the global state and outputs a joint action for all intersections.
Although theoretically capable of finding globally optimal policies, this approach
is computationally intractable in large networks because the joint action space
grows exponentially with the number of agents. \textit{Decentralised} methods allow
each agent to learn and act independently from local observations, scaling well to
large networks but introducing \textit{non-stationarity}: from any single agent's
perspective, the environment appears non-stationary because the other agents are
simultaneously updating their policies.

\subsection{Independent Learning with Parameter Sharing}

This thesis adopts \textit{Independent Learning with Parameter Sharing}: each agent
observes and acts locally at its own intersection, but all agents share a single
neural network~\cite{terry2020revisiting}. This architecture is applied to all three
algorithms, yielding the variants IPPO, IDQN, and IDDQN. Parameter sharing brings
two practical benefits: it reduces the total number of trainable parameters, and it
allows experience collected at every intersection to contribute to a single
centralised gradient update, substantially improving data efficiency. Despite its
simplicity, this approach has been shown to match or surpass more complex
centralised methods in competitive multi-agent
benchmarks~\cite{dewitt2020independent}.

% ================================================================
\section{Simulation Environment and Dataset}
% ================================================================

\subsection{SUMO Traffic Simulator}

All experiments are conducted in \textit{SUMO}, an
open-source, microscopic, multi-modal traffic simulator developed by the German
Aerospace Center (DLR)~\cite{lopez2018sumo}. SUMO models individual vehicle
behaviour---including car-following, lane-changing, and intersection management---at
a sub-second resolution, making it widely used in RL-based traffic signal control
research~\cite{marl_tsc_sustainability2023}. Agent--environment interaction is
handled through the \textit{TraCI} API, which exposes
per-lane vehicle counts, queue lengths, waiting times, and current signal phases in
real time, and allows external controllers to override signal timings at every
simulation step.

\subsection{TAPAS Cologne Dataset}

The traffic demand used in both experimental scenarios is derived from the
\textit{TAPAS Cologne} dataset~\cite{tapas_cologne_sumo}, which describes the
complete vehicle trip demand within the city of Cologne, Germany, over a 24-hour
period. The demand is generated by the TAPAS activity-based model, which synthesises
individual travel decisions based on German mobility surveys and detailed land-use
data. With approximately 1.2 million vehicle trips distributed across a full-day
simulation, TAPAS Cologne is recognised as one of the largest freely available
traffic simulation datasets and provides realistic, time-varying congestion patterns
that go well beyond synthetic benchmarks.

Two sub-scenarios derived from this dataset are used in this thesis: \textbf{Cologne-3},
a single-axis arterial corridor, and \textbf{Cologne-8}, a $2{\times}2$ signalised
grid network. Both scenarios are run at 30\% of the full demand scale---the
recommended setting in the official SUMO documentation to avoid network-wide
gridlock---while still preserving realistic peak-hour congestion dynamics.